\section{Financial Economics Insights}
\label{sec:econ}

The preceding experiments demonstrate that DeXposure-FM can forecast network structure and detect stress propagation patterns. This section discusses how these capabilities translate into actionable tools for financial oversight, the conditions under which model outputs remain reliable, and practical considerations for policymakers.

\subsection{Macroprudential Applications}
\label{subsec:macro-applications}

DeXposure-FM generates several outputs directly relevant to macroprudential monitoring of decentralised financial infrastructure, addressing the data gaps highlighted by international standard-setters \citep{FSB2023,IOSCO2023DeFi}.

\paragraph{Systemic importance ranking.}
The protocol-level Systemic Importance Score (SIS) defined in Equation~\eqref{eq:sis} combines network centrality (PageRank), exposure concentration (TailExposure), and scale (TVL) into a single metric. This approach draws on the network-theoretic insight that interconnectedness, not just size, determines systemic fragility \citep{Acemoglu2015SystemicRisk,GlassermanYoung2016Contagion}. Regulators can use weekly SIS rankings to identify protocols whose failure would most disrupt the broader ecosystem. Unlike traditional banking metrics that rely on quarterly balance-sheet filings, SIS updates weekly from on-chain data, enabling more timely surveillance.

\paragraph{Sector spillover monitoring.}
The sector-to-sector spillover matrix $\mathbf{S}$ (Section~\ref{subsubsec:systemic-risk}) reveals which protocol categories are net exporters or importers of risk, analogous to the directional spillover measures proposed by \cite{DieboldYilmaz2012} for traditional asset markets. Our empirical results indicate that stablecoins and bridges frequently occupy central positions in cross-sector flows, consistent with their role as liquidity conduits \citep{KosseEtAl2023Stablecoin,ESRB2025}. Tracking changes in off-diagonal elements of $\mathbf{S}$ over time can alert supervisors to emerging concentration.

\paragraph{Stress testing and contagion assessment.}
The DebtRank-style contagion simulation (Section~\ref{subsubsec:contagion}) allows regulators to pose ``what-if'' questions, following the loss-propagation logic of \cite{EisenbergNoe2001} and \cite{BattistonEtAl2012DebtRank}. Our simulations on 2025 network snapshots suggest that isolated shocks propagate minimally (system loss $<$0.1\%), reflecting diversified counterparty structures. However, coordinated shocks to bridge protocols produce deeper contagion (up to 0.2\% system loss across 12 affected protocols), highlighting bridges as potential amplification nodes---a concern echoed by recent policy analyses \citep{Aufiero2025}.

\paragraph{Early warning signals.}
Rapid changes in network density ($\Delta\rho$), concentration (HHI spikes exceeding two standard deviations), and degree assortativity can precede stress events. Classical financial-network theory predicts that denser networks may initially diversify risk but become channels for contagion beyond a critical threshold \citep{Acemoglu2013,GaiKapadia2010Contagion}. Our shock-event analysis (Table~\ref{tab:shock-results}) shows that the Terra/Luna collapse was accompanied by a 33.5\% TVL contraction and a 9.4\% increase in edge count as protocols restructured exposures. Integrating such structural indicators into a monitoring dashboard could provide regulators with lead time to investigate anomalies.


\subsection{Validity and Scope}
\label{subsec:validity}

Model-based risk measures are only as reliable as the conditions under which they were trained and validated \citep{MullainathanSpiess2017}. We identify several boundaries that users should keep in mind.

\paragraph{Stable versus turbulent regimes.}
The DeXposure dataset exhibits 98.5\% mean edge overlap week-to-week, meaning most credit relationships persist. Under such stable conditions, GraphPFN achieves AUROC $>$0.91 across all forecast horizons (Table~\ref{tab:main-results}). During structural breaks---new protocol launches, governance attacks, or rapid deleveraging---the historical edge distribution may shift abruptly, and forecasts should be interpreted with greater caution. The degradation of ROLAND at $h=14$ (AUROC 0.718) illustrates how models without pre-trained representations struggle when asked to extrapolate beyond short horizons.

\paragraph{Coverage.}
DeXposure-FM is trained exclusively on on-chain DeFi protocols sourced from DeFiLlama. It does not observe:
\begin{itemize}
  \item Centralised exchange (CEX) balances or order-book exposures,
  \item Over-the-counter (OTC) derivatives or lending agreements,
  \item Off-chain reserves backing stablecoins (e.g., Treasury holdings of USDC) \citep{AnaduEtAl2023StablecoinsMMF}.
\end{itemize}
Consequently, the model captures intra-DeFi contagion but cannot directly assess spillovers to or from traditional financial institutions---a gap that matters as DeFi-TradFi linkages deepen \citep{ESRB2025}.

\paragraph{Temporal resolution.}
The dataset aggregates to weekly snapshots. Intraday or even daily dynamics---such as flash-loan attacks \citep{Qin2021FlashLoans} or rapid liquidation cascades \citep{Warmuz2023ToxicLiquidationSpirals}---are smoothed out. For high-frequency risk monitoring, complementary tools operating on block-level data would be required.


\subsection{Practical Caveats for Policymakers}
\label{subsec:caveats}

\paragraph{Complement, not replace, qualitative judgment.}
Model outputs such as SIS rankings and spillover indices are quantitative summaries, not definitive verdicts. A protocol may rank low on SIS yet pose idiosyncratic risks (e.g., smart-contract vulnerabilities, oracle manipulation) that network topology alone cannot capture \citep{DuleyEtAl2023Oracle,Gogol2024SoKDeFi}. Supervisors should triangulate model signals with code audits, governance analyses, and market intelligence.

\paragraph{TVL as a proxy for exposure.}
Total Value Locked is a convenient but imperfect measure \citep{Bertomeu2024}. It aggregates heterogeneous assets---stablecoins, volatile tokens, LP positions---without adjusting for liquidity or collateralisation quality. A protocol with \$1B in illiquid governance tokens is not equivalent to one with \$1B in USDC. Future work could incorporate asset-level risk weights to refine exposure estimates.

\paragraph{Model drift and retraining.}
DeFi evolves rapidly: new chains launch, token standards change, and incentive mechanisms shift. A model trained on 2020--2024 data may gradually lose accuracy as the ecosystem diverges from historical patterns. Periodic retraining on updated snapshots is essential to maintain forecast quality.

\paragraph{Cross-chain data consistency.}
DeXposure aggregates data across 602 blockchains, each with its own indexing quirks and oracle delays. Minor inconsistencies in timestamp alignment or token-price feeds can introduce noise. Users should be aware that edge weights near reporting boundaries may reflect data artefacts rather than genuine exposure changes.
